\inicializarSubseccion{Conclusión de Andres}

\tab En el reto aprendí muchísimas cosas sobre lógica proposicional y sobre los usos que tienen los polinomios en la creación de objetos. Es importante saber que previo a este reto, yo ya tenía cierto conocimiento sobre la interpolación de puntos para crear un objeto, pero nunca lo había puesto en práctica. Aprendí a manejar mejor MATLAB, la cúal me permitió entender bastante teoría para realizar un acercamiento a la solución final del reto. Agregado a esto, también se le puede atribuir que aprendí varios métodos para el manejo de nuestros arrays de puntos (por ejemplo, el uso de la función `vander` para poder obtener los coeficientes de mis funciones). En definitiva, mi mayor aprendizaje estuvo en el área de la programación, ya que aprendí a utilizar muchísima funcionalidad clave en el lenguaje de programación de excelencia para ingenierías.

\tab Diría que las cosas que más se me complicaron fue la corrección del polinomio y la visualización 3D. Más que nada fue la simple razón que el método que fue utilizado para el reto no fue el más adecuado como para realizar un producto. Existen técnicas de conexiones de puntos que puedes poner $n$ puntos y te genera un polinomio de $n-1$ puntos. Esto es importante ya que a la hora de corregir u optimizar material para la figura que estás creando, es más sencillo modificar una sola función en lugar de varias funciones independientes. La visualización en 3D fue un tema nos dio problemas ya que el documento anexado mostraba un método que funciona muy bien para proyectar un objeto, pero se rompe inmediatamente al agregar otro objeto. Esto hizo que tuviéramos que buscar otros métodos para poder mostrar las figuras de manera correcta en nuestro entregable final.

\tab Para finalizar, es importante mencionar que los cálculos jugaron un papel importante, sin embargo, diría yo que no fue el más crítico para este proyecto. Debido a la manera en la que interpolamos nuestros bocetos e imágenes, casi toda la matemática compleja del proyecto no jugó parte en el reto. Sin embargo, si no fuera por estos principios de modelación matemática, este reto definitivamente no se hubiera materializado de la manera en la que se hizo; recalcando aquí el uso de la teoría para poder hacer proyectos realidad.
